\section{Comparación y relación entre los dos modelos}

\subsection{VAE vs. GAN: comparación de arquitecturas}

\begin{figure}[H]
\centering
\includegraphics[width=0.85\textwidth]{slides-images/slide-46.jpg}
\caption{Dos grandes familias de modelos de variable latente: VAE y GAN\protect\footnotemark}
\end{figure}
\footnotetext{Diapositiva 46.}

\begin{table}[H]
\centering
\caption{Comparación sistemática entre VAE y GAN}
\begin{tabular}{lll}
\toprule
\textbf{Dimensión} & \textbf{VAE} & \textbf{GAN} \\
\midrule
Método de modelado & Modela explícitamente la densidad $p(x)$ & Aprende implícitamente el proceso de muestreo \\
Estructura de red & Codificador + decodificador & Generador + discriminador \\
Método de entrenamiento & Maximiza el ELBO & Juego adversario minimax \\
Espacio latente & Generado por el codificador, estructurado & Espacio de ruido, sin codificador \\
Calidad de generación & Generalmente más borrosa & Generalmente más nítida y definida \\
Estabilidad del entrenamiento & Relativamente estable & Inestable, propenso al colapso de modo \\
Interpretabilidad & Espacio latente interpretable & Espacio latente con interpretabilidad débil \\
\bottomrule
\end{tabular}
\end{table}

\begin{warningbox}{Por qué el VAE genera imágenes borrosas}
El VAE utiliza una pérdida de reconstrucción a nivel de píxel (como el MSE), lo cual hace que el modelo tienda a producir el «promedio» de varios resultados posibles, generando salidas borrosas. El GAN, en cambio, optimiza la calidad de generación mediante la señal adversaria que proporciona el discriminador, sin esta limitación, por lo que suele generar imágenes más nítidas. Sin embargo, al GAN le falta un codificador explícito como el del VAE, lo que dificulta la inferencia y la manipulación del espacio latente.
\end{warningbox}

\subsection{Resumen de la sección}

El VAE y el GAN son dos paradigmas complementarios de modelos generativos. El VAE ofrece un espacio latente estructurado y un entrenamiento estable, pero su calidad de generación es limitada; el GAN ofrece una calidad de generación sobresaliente, pero su entrenamiento es inestable y carece de un espacio latente explícito. En aplicaciones prácticas, se puede elegir el modelo adecuado según las necesidades, o usar arquitecturas híbridas que combinen las ventajas de ambos.

